\section{결론}
\label{sec:conclusion}

\begin{revision}

비전--언어 주행 모델의 CoC는 행동의 이유를 설명하지만, 그 행동 요구가 언제까지 유지되고 무엇이 충족되면 끝나는지는 대개 밝히지 않는다. 그래서 각각 자연스러운 설명도 시간순으로 이어 읽으면 앞선 의무와 충돌하거나 해제 근거 없이 진행을 정당화할 수 있으며, 장면 경계를 넘으면 같은 대상에 관한 의무인지조차 확인하기 어렵다. 이런 주석을 학습 자료나 검토 근거로 그대로 쓰면 행동 요구가 모호해지고, 같은 자료를 다시 검사할 때 어떤 조건과 근거를 적용해야 하는지도 불분명해진다. 따라서 조건부 행동이 이어지는 CoC를 학습 자료나 검토 근거로 쓰려면 앞선 의무와 해제 근거의 연결을 확인하는 검사가 필요하다.

이를 위해 본 논문은 연속 CoC에서 의무를 해제 조건 및 조건별 증거와 연결하고, 사건 순서에 따라 확정 모순과 근거 부족을 구분하며, 최초 문제 사건과 관련 의무를 해당 원문 구절과 함께 제시하는 검사 프레임워크를 제안했다. 이전 상태를 기억하고 판정을 참ㆍ거짓ㆍ미결정의 세 값으로 내는 방식은 기존 실행 중 검증의 기법이며, 본 논문의 기여는 이를 자동 생성 주석에 적용하기 위한 입력 표현과 판정 규칙에 있다.

작성된 미해제 충돌 사례에서는 이전 의무의 보존이 해당 사건 관계를 판정하는 데 필요했다. 다만 이 관찰을 모든 CoC에 대한 필요성ㆍ효과나 UPPAAL의 정확도 우위로 일반화하지 않는다. 또한 기억만으로는 충분하지 않았다. 전체 이력을 기억하는 기존 검사기도 144건 중 96건만 기대 판정과 일치했으며, 실패는 여러 의무를 하나의 채널에 담고 대기 중 미상을 누적하는 표현과 정책에서 생겼다. 의무별 상태와 조건별 증거를 나눈 조건 연결 모델은 같은 사례에서 판정과 문제 위치를 모두 기대대로 돌려주었으며, 이는 개발 후 동일 사례 재시험의 결과다. 단순 FSM도 같은 판정을 냈으므로 UPPAAL의 의의는 정확도가 아니라 판정 규칙과 검사 속성을 명시하고 진단 경로를 남기는 데 있다.

장면 경계에서는 대상 동일성, 적용 범위, 증거 유효성이 주어질 때 상태 유지가 판정에 필요했다. 경계에서 의무를 버리는 모델은 확정 모순 44건 중 41건을 양립으로 판정했다. 조건 연결 모델과 경계 확장 모델의 상태 갱신 구현은 선별 없이 열거한 3,500만여 개의 사건열에서 각각 대응하는 수식 판정과 일치했다.

실제 CoC에서는 사람의 판정부터 크게 엇갈렸고(텍스트 일관성 $\kappa=0.027$), 독립된 정답이 없어 오류 검출 성능을 산출하지 않았다. 현재 결과는 명시된 프레임의 시간ㆍ조건 관계를 실행하고 검토 근거를 추적하는 방법 실증이다. 자동 의미 추출, 외부 작성 시험, 자연 CoC의 의미 정답 구축과 LLM 판정기와의 비교는 \ref{sec:future-work}절의 후속 과제로 남는다.
\end{revision}
